%%%%%%%%%%%%%%%%%%%%%%%%%%%%%%%%%%%%%%%%%%%%%%%%%%%%%%%%%%%%%%%%%%%%%%%%%%%%%%%%%%%%%%%%%%%%%%%%%%%%
\chapter{Related Work} %%%%%%%%%%%%%%%%%%%%%%%%%%%%%%%%%%%%%%%%%%%%%%%%%%%%%%%%%%%%%%%%%%%%%%%%%%%%%
%%%%%%%%%%%%%%%%%%%%%%%%%%%%%%%%%%%%%%%%%%%%%%%%%%%%%%%%%%%%%%%%%%%%%%%%%%%%%%%%%%%%%%%%%%%%%%%%%%%%

Understanding the challenges novice programmers face when dealing with compiler error messages is a topic of growing importance in computing education research.
Compiler diagnostics are often the first line of feedback that learners receive during programming tasks.
However, this feedback is frequently described as inadequate, cryptic, and even detrimental to learning.
This chapter synthesizes findings from two key studies in this area:
Becker et al.'s extensive literature review on compiler error messages \cite{becker},
and Barik et al.'s empirical eye-tracking study that explores how developers interact with such messages in practice \cite{barik}.\\[5pt]
\noindent Becker et al. \cite{becker} conducted a comprehensive review of over 300 publications spanning several decades,
categorizing research according to pedagogical relevance, technical challenge, enhancement efforts, and empirical findings.
Their work emphasizes the longstanding and pervasive problem of unhelpful error messages, particularly in introductory programming contexts.
Barik et al. \cite{barik} complement this with rigorous experimental data:
their eye-tracking study with 56 novice and intermediate participants reveals not only that developers do read compiler error messages,
but also that these messages demand cognitive effort comparable to reading source code, and that difficulty in interpreting them significantly predicts task performance.\\[5pt]
\noindent These two studies provide a robust framework for understanding both the qualitative and quantitative aspects of compiler error messages,
and form the foundation for discussing their educational implications and potential improvements.

%%%%%%%%%%%%%%%%%%%%%%%%%%%%%%%%%%%%%%%%%%%%%%%%%%%%%%%%%%%%%%%%%%%%%%%%%%%
\section{Challenges for Novice Programmers} %%%%%%%%%%%%%%%%%%%%%%%%%%%%%%%
%%%%%%%%%%%%%%%%%%%%%%%%%%%%%%%%%%%%%%%%%%%%%%%%%%%%%%%%%%%%%%%%%%%%%%%%%%%

Learning to program involves mastering syntax, logic, and problem decomposition—skills that are cognitively demanding on their own.
Compiler error messages, instead of supporting the learning process, often act as additional hurdles.
Becker et al. \cite{becker} stress that from the early days of computing education, these messages have been described as inadequate, confusing, and unhelpful, especially for beginners.
Their review cites evidence as early as the 1960s and includes anecdotes from instructors and students alike, such as one student claiming: ``If compiler errors were worth one Euro, I would be a millionaire.''\\[5pt]
\noindent Becker et al. identify two categories of difficulty:

\begin{itemize}
    \item \textbf{Internal difficulties:} Emotional reactions such as frustration, anxiety, and intimidation are frequently triggered by error messages.
        These reactions reduce learners' confidence and self-efficacy, leading to lower persistence.
    \item \textbf{External difficulties:} Error messages often contain terminology that assumes a level of programming expertise.
        Moreover, the feedback may be imprecise or even misleading—such as one error producing multiple, contradictory messages depending on compiler or context.
\end{itemize}

\noindent Barik et al. \cite{barik} provide empirical reinforcement for these observations.
Their eye-tracking data revealed that:

\begin{itemize}
    \item Developers allocate between 13-25\% of their total task time to reading compiler messages.
        For a four-hour assignment, this amounts to roughly 30-60 minutes focused solely on error interpretation.
    \item Reading difficulty of error messages was statistically significant in predicting the correctness of solutions (\(p < .0001\)), with an \(R^2\) value of \(0.16\), indicating a non-negligible effect.
    \item Most participants repeatedly revisited the same error message areas, suggesting cognitive overload and uncertainty about how to proceed.
\end{itemize}

\noindent Barik et al. also introduced the concept of ``revisits'' as a proxy for cognitive difficulty, showing that increased revisits to error messages correlated with lower task performance.
This aligns with earlier findings from psychology that repeated eye fixations are markers of comprehension difficulty.\\[5pt]
\noindent Together, these studies show that novice programmers are not only slowed down by error messages,
but their overall learning trajectories may be compromised when the messages are ambiguous, misleading, or too technical.

%%%%%%%%%%%%%%%%%%%%%%%%%%%%%%%%%%%%%%%%%%%%%%%%%%%%%%%%%%%%%%%%%%%%%%%%%%%
\section{Role of Error Messages in Learning} %%%%%%%%%%%%%%%%%%%%%%%%%%%%%%
%%%%%%%%%%%%%%%%%%%%%%%%%%%%%%%%%%%%%%%%%%%%%%%%%%%%%%%%%%%%%%%%%%%%%%%%%%%

Compiler error messages serve as one of the few immediate feedback mechanisms available to learners when an instructor is not present.
From a pedagogical standpoint, these messages could be powerful tools for formative assessment—guiding learners to identify misconceptions, reflect on their understanding, and improve incrementally.
However, Becker et al. \cite{becker} and Barik et al. \cite{barik} highlight a gap between this potential and the reality of current message design.\\[5pt]
\noindent Barik et al. categorize the factors influencing student engagement with error messages into internal and external factors:

\begin{itemize}
    \item \textbf{Internal factors:} Many learners find error messages overwhelming or accusatory in tone.
        Such emotional responses are not trivial; they directly impact engagement and willingness to persist in debugging tasks.
        Some learners internalize the messages as a reflection of their own inadequacy, which can reinforce imposter syndrome.
    \item \textbf{External factors:} Learning environments often do not support constructive interaction with error messages.
        The design of many IDEs prioritizes functionality over usability, and educational tools often lack scaffolding features like message disambiguation,
        inline explanations, or novice-friendly vocabulary.
\end{itemize}

\noindent Becker et al. argue that compiler error messages are often misaligned with the kinds of mistakes students actually make.
The error messages may accurately describe what the compiler failed to parse or infer, but not why the student made the mistake in the first place.
This disconnect can lead students to apply ineffective or even counterproductive debugging strategies, such as removing entire code blocks until the error disappears.\\[5pt]
\noindent Moreover, learners often resort to \textit{externalizing} the debugging process—copy-pasting errors into search engines or
generative AI models without developing a meaningful understanding of the underlying issue.
While such strategies may offer short-term fixes, they hinder long-term skill acquisition.

%%%%%%%%%%%%%%%%%%%%%%%%%%%%%%%%%%%%%%%%%%%%%%%%%%%%%%%%%%%%%%%%%%%%%%%%%%%
\section{Existing Approaches to Improving Error Messages} %%%%%%%%%%%%%%%%%
%%%%%%%%%%%%%%%%%%%%%%%%%%%%%%%%%%%%%%%%%%%%%%%%%%%%%%%%%%%%%%%%%%%%%%%%%%%

Various approaches have been proposed and implemented to improve the usability of error messages, especially for beginners.
Becker et al. classify these approaches into three major categories:

\begin{itemize}
    \item \textbf{Enhanced Error Messages:} These are efforts to rewrite or augment the default messages produced by compilers.
        Enhancements include simplifying technical language, offering concrete examples, and suggesting specific code changes.
    \item \textbf{Intelligent Feedback Systems:} Educational IDEs and learning platforms have begun to integrate features such as syntax-aware hint systems, Quick Fix options, and context-sensitive help.
        Some tools also utilize AI-based tutoring systems to adapt feedback based on user history.
    \item \textbf{Community and Data-Driven Solutions:} Platforms like Stack Overflow, along with repositories of commonly asked questions and typical fixes,
        help crowdsource explanations and debugging strategies. Some environments are beginning to embed such resources directly into the IDE.
\end{itemize}

\noindent Despite these promising developments, Becker et al. caution that many of these systems are still reactive and heuristic in nature.
For instance, Quick Fix suggestions may resolve a syntax issue but reinforce a faulty conceptual model.
Moreover, compiler behaviors are inconsistent: the same logical error might yield different messages across compilers or languages, or even different versions of the same compiler.
This inconsistency further complicates error interpretation for learners.

%%%%%%%%%%%%%%%%%%%%%%%%%%%%%%%%%%%%%%%%%%%%%%%%%%%%%%%%%%%%%%%%%%%%%%%%%%%
\section{Educational Feedback Theories} %%%%%%%%%%%%%%%%%%%%%%%%%%%%%%%%%%%
%%%%%%%%%%%%%%%%%%%%%%%%%%%%%%%%%%%%%%%%%%%%%%%%%%%%%%%%%%%%%%%%%%%%%%%%%%%

The effectiveness of error messages should be evaluated through the lens of educational feedback theory.
According to Becker et al. \cite{becker}, error messages can be thought of as a form of formative feedback, which is essential for learning when it is:

\begin{itemize}
    \item \textbf{Timely:} Delivered immediately after the error occurs.
    \item \textbf{Specific:} Clearly identifies the nature and location of the error.
    \item \textbf{Actionable:} Provides guidance on how to resolve the issue and avoid it in the future.
\end{itemize}

\noindent In addition, feedback should support a \textit{growth mindset} by framing errors as learning opportunities.
Becker et al. argue that current compiler messages often fall short in this regard. Instead of promoting understanding, they frequently confuse or discourage learners.
Some studies advocate for adaptive feedback systems that adjust based on the user's proficiency.\\[5pt]
\noindent For example, beginners could receive verbose messages with explanations, while advanced users might prefer concise diagnostics.
This personalization of feedback is consistent with Vygotsky's concept of the Zone of Proximal Development (ZPD),
where instructional scaffolding enables learners to tackle challenges just beyond their current capability.\\[5pt]
\noindent Barik et al. suggest that IDEs and compilers could benefit from being ``situationally aware''—adapting their diagnostics based on coding context, user behavior, and even gaze tracking data.
Such intelligent systems would align with principles of human-computer interaction and could lead to more meaningful and effective educational interventions.

%%%%%%%%%%%%%%%%%%%%%%%%%%%%%%%%%%%%%%%%%%%%%%%%%%%%%%%%%%%%%%%%%%%%%%%%%%%
% References %%%%%%%%%%%%%%%%%%%%%%%%%%%%%%%%%%%%%%%%%%%%%%%%%%%%%%%%%%%%%%
%%%%%%%%%%%%%%%%%%%%%%%%%%%%%%%%%%%%%%%%%%%%%%%%%%%%%%%%%%%%%%%%%%%%%%%%%%%

% Note: Ensure these entries exist in your BibTeX file.
% \bibitem{becker} Becker et al. (2019). Compiler Error Messages Considered Unhelpful: The Landscape of Text-Based Programming Error Message Research.
% \bibitem{barik} Barik et al. (2017). Do Developers Read Compiler Error Messages?

